\clearpage
\scanpage{86}
\mspage{77}
% archive sheet number: 86

\begin{samepage}
satisfait la propriété suivante
\[
\forall F\in\unsure{\mathscr F|\mathcal F},\quad\text{on a}\quad
\unsure{\varphi|\phi}(\operatorname{cosupp}^{\circ}(F))=\tikz[remember picture,baseline=(s86xone.base)]{\node[draw,circle,inner sep=1pt] (s86xone) {$X$};}\setminus|F|
\]
\reversemarginpar
\margin{\tikz[remember picture]{\node[inner sep=0pt] (s86note) {$?$ \uline{plutôt} $P$};}}
\normalmarginpar
% The marginal correction has arrows to the first two circled X's and the X after the quantifier in the boxed line.
i.e.
\[
\begin{aligned}
\tikz[remember picture,baseline=(s86xtwo.base)]{\node[draw,circle,inner sep=1pt] (s86xtwo) {$X$};}\setminus|F|&=\bigcup_{\substack{Z\in M\\\text{t.q. }\tikz[remember picture,baseline=(s86cond.base)]{\node[inner sep=0pt] (s86cond) {$\scriptstyle Z\mathrel{|\circ|}F$};}}}|Z|^{\circ}\\
&\bigl[\text{i.e. }|Z|^{\circ}\subset X\setminus|F|\bigr]
\end{aligned}
\]
\margin{\tikz[remember picture]{\node[text width=\marginparwidth,inner sep=0pt] (s86right) {\deleted{Sous cette forme, la condition implique la fidélité}};}}
% The cancelled right-margin note has a left-pointing arrow to the union condition.
i.e.
\[
\boxed{\forall x\in \tikz[remember picture,baseline=(s86xthree.base)]{\node[inner sep=0pt] (s86xthree) {$X$};}-|F|,\quad\exists Z\text{ avec }x\in|Z|^{\circ}\subset X\setminus|F|}
\]
\figure{\begin{tikzpicture}[remember picture,overlay]
\draw[->,gray] (s86note.east) to[out=15,in=-120] (s86xone.south west);
\draw[->] (s86note.east) to[out=25,in=-110] (s86xtwo.south west);
\draw[->] (s86note.east) to[out=0,in=90] (s86xthree.north);
\draw[->] (s86right.west) to[out=190,in=0] (s86cond.east);
\end{tikzpicture}}
\end{samepage}
Ceci posé, on trouve

\uline{Proposition} Soit $\unsure{\Phi|\mathscr P}\subset F$. Alors
\[
\varphi(\operatorname{supp}^{\circ}(\Phi))=|\Phi|^{\circ}
\left(\overset{\mathrm{def}}{=}\bigcup_{X\in\Phi}|X|^{\circ}\right)
\tag*{(*)}
\]

\uuline{Dém} Comme $\deleted{\unsure{X\in}}\Phi\subset\operatorname{supp}^{\circ}(\Phi)$, on a $\deleted{\varphi(\Phi)\subset}\varphi(\Phi)\overset{\mathrm{def}}{=}|\Phi|^{\circ}\subset\varphi(\operatorname{supp}^{\circ}(\Phi))$, i.e. on a $\supset$. Prouvons $\subset$, \deleted{Soit $Y\in\operatorname{supp}^{\circ}\Phi$, i.e.}
% slanted
\margin{expliciter d'abord la relation $\Phi\longrightarrow|\Phi|^{\circ}$, en termes de la figure image de $F$ dans $\unsure{P|\mathscr P}$\ldots}
\[
\forall Y\in\operatorname{supp}^{\circ}(\Phi)\Longrightarrow |Y|^{\circ}\subset|\Phi|^{\circ}\qquad\text{i.e.}
\]
\[
\begin{aligned}
|Y|^{\circ}\not\subset|\Phi|^{\circ}&\Longrightarrow Y\notin\operatorname{supp}^{\circ}(\Phi),\quad\text{i.e. }\exists Z\text{ avec}\\
\text{i.e. }|Y|^{\circ}\cap\complement|\Phi|^{\circ}\neq\varnothing
&\qquad Z\mathrel{|\circ|}\Phi,\quad\text{et }Z\mathrel{\overline{|\circ|}}Y,\quad\text{i.e.}\\
&|Z|^{\circ}\cap|\Phi|^{\circ}=\varnothing,\quad|Z|^{\circ}\cap|Y|^{\circ}\neq\varnothing.
\end{aligned}
\]
Distinguons deux cas

cas $|Y|^{\circ}\cap\bigl(\underbrace{|F|\setminus|\Phi|^{\circ}}_{=|F\setminus\Phi|^{\circ}}\bigr)\neq\varnothing$, \deleted{\unsure{Si $x$ est dans le}}
\[
\exists Z\in F\setminus\Phi,\quad\text{tel que }|Y|^{\circ}\cap|Z|^{\circ}\neq\varnothing.\quad\text{OK}
\]
cas $|Y|^{\circ}\cap(P\setminus|F|)\neq\varnothing$, alors on applique l'hypothèse de pleine fidélité à $F$.
